\subsection{Inverses with Even Powers} 
When we want to take the inverse of an even root, we need to be careful of the \makeblank[1in] of the resulting function.
\begin{Example}Let $f(x)=\sqrt{x+3}+4$. Find $f\inv$.
\begin{lpic}{}{8}
\foreach \y/\lbl in {1/Swap $x$ and $y$,
										2/Isolate the root,
										3/Consider the domain,
										5/Rewrite as a power function,
										7/Finish solving for $y$,
										8/Rewrite in function notation}
	\regtextleft{0.25}{-\y}{\lbl};
\end{lpic}
\end{Example}
\noi We can't take an inverse of an even power function unless the \makeblank[1in] is restricted. If it is, we need to be careful of whether we want the \makeblank[1in] or \makeblank[1in] root, so the range of the inverse matches the domain restriction. (We can't use both if we want to get a function!)
\begin{Example}Let $g(x)=(2-x)^4-7,\ x\geq 2$. Find $g\inv$.
\begin{lpic}{}{10}
\foreach \y/\lbl in {1/Swap $x$ and $y$,
										2/Isolate the 4th power,
										3/Consider the range,
										6/Rewrite as a radical function,
										8/Finish solving for $y$,
										10/Rewrite in function notation}
	\regtextleft{0.25}{-\y}{\lbl};
\end{lpic}
\end{Example}
